% Chapter — Built-in reference GNC components (FR-25). Follows the
% mandatory FR-29 six-section template: purpose; assumptions and domain of
% validity; derivation; discretization and implementation; validation
% evidence; references. Written ahead of the implementation (Phase 6):
% this chapter is the binding specification for the built-in
% IGncComponent implementations and the plugin-boundary semantics they
% share. The implementing change adds the chapter-manifest entry and the
% \testid{} evidence table. The equation labels eq:gnc:deltaq, eq:gnc:pd,
% eq:gnc:sat, eq:gnc:werr, eq:gnc:exactrot, eq:gnc:drprop, and
% eq:gnc:cmdrate are intended to be echoed verbatim in the module source
% (FR-29 traceability).
\chapter{Built-in Reference GNC Components}
\label{ch:gnc-builtin}

\section{Purpose}
\label{sec:gnc:purpose}

This chapter specifies the built-in C++ reference components behind the
FR-25 plugin interface: (a) a quaternion-error PD attitude controller,
(b) a dead-reckoning attitude navigator driven by the
Chapter~\ref{ch:sensors-imu} gyro increments, and (c) the ascent
pitch-program guidance --- the Phase~4 open-loop pitch program of
Chapter~\ref{ch:vehicle6dof} re-expressed as an \texttt{IGncComponent}
so the same ascent flies under the GNC stack (Phase~6 exit
criterion~10). It also pins the plugin-boundary semantics every
component, built-in or user-supplied, shares: zero-order-hold commands
over the major cycle (D-5), the \texttt{latency\_cycles} shift (exit
criterion~8), and the \texttt{oracle} truth-injection flag (exit
criterion~5). These components are \emph{reference} implementations:
deliberately simple, exactly specified, and reproducible in a few lines
of Python --- exit criterion~2 requires a Python reimplementation of the
PD controller to match the C++ torques to
$< \qty{e-9}{\newton\metre}$, so every branch and operation order in
Section~\ref{sec:gnc:pd} is normative.

\section{Assumptions and domain of validity}
\label{sec:gnc:domain}

Assumptions:
\begin{enumerate}
  \item \textbf{Cycle-synchronous execution.} Components run once per
    major cycle on the D-5 schedule; all inputs are the sensor samples
    of that cycle and all outputs are held (zero-order hold) across the
    next application interval. No component sees intra-cycle states.
  \item \textbf{No truth in \texttt{GncInput}.} Components consume only
    sensor outputs, their own state, and configuration. The
    scenario-level \texttt{oracle: true} debug flag injects truth into
    \texttt{GncInput} \emph{and} stamps the log header, so an oracle
    run is identifiable from its header alone (exit criterion~5); the
    reference components themselves never read the oracle fields.
  \item \textbf{Diagonal gains, per-axis saturation.} The PD controller
    uses positive per-axis gain vectors and clamps each torque axis
    independently (equation~\eqref{eq:gnc:sat}); a saturated command is
    direction-distorted, which is the standard, simplest behavior and
    is stated rather than hidden.
  \item \textbf{Dead reckoning is uncompensated.} The dead-reckoning
    navigator applies no bias compensation and no coning correction;
    its drift is exactly the integrated gyro error plus the bounded
    single-sample non-commutativity residual of
    equation~\eqref{eq:gnc:coningbound}. It exists as the minimal
    reference navigator and as the mechanization core the
    Chapter~\ref{ch:ekf} filter shares.
  \item \textbf{Guidance restates Phase 4 math.} The ascent guidance
    reuses the pitch-program geometry of Chapter~\ref{ch:vehicle6dof}
    (equations~\eqref{eq:vehicle6dof:pitchaxis}
    and~\eqref{eq:vehicle6dof:attitude}) unchanged; this chapter adds
    only its packaging as a component and the commanded-rate
    convention.
\end{enumerate}

Domain of validity: attitude-regulation regimes with per-cycle error
rotations far from the $\pm\pi$ antipode (the sign branch of
equation~\eqref{eq:gnc:pd} handles the double cover; exactly at
$\delta q_w = 0$, $180^{\circ}$ error, the branch picks $s = +1$
deterministically and the controller still converges, though the
commanded direction at that single point is convention, not physics).
Gains are the user's stability responsibility: the PD law is globally
stabilizing for a rigid body with pure torque actuation and positive
gains (Wie~\cite{wie2008}), but actuator saturation, latency, and
flexible dynamics can destabilize aggressive tunings --- the simulator
reports what the configuration does, per the project's
no-silent-correction rule (DX-2).

\section{Derivation}
\label{sec:gnc:derivation}

\subsection{Plugin-boundary semantics}
\label{sec:gnc:interface}

Let $T_c$ be the control period and $t_k = k\,T_c$ the cycle
boundaries. At cycle $k$ the scheduler samples sensors, calls each
configured component's \texttt{update}, composes the stack's final
actuator command vector $\vv{u}_k$ (torque and/or gimbal/throttle
commands, per the actuator suite of Chapter~\ref{ch:actuators}), and
applies it as a zero-order hold. With the configured latency $L =
\texttt{latency\_cycles} \ge 0$, the command \emph{applied} over
$[t_k, t_{k+1})$ is
\begin{equation}
  \vv{u}^{\mathrm{applied}}_k =
  \begin{cases}
    \vv{u}_{k-L}, & k \ge L,\\[2pt]
    \vv{u}_{\mathrm{neutral}}, & k < L,
  \end{cases}
  \label{eq:gnc:latency}
\end{equation}
where the neutral command is zero torque, zero gimbal deflection, and
throttle hold at the sequence-commanded value. The shift applies once,
at the GNC-to-plant boundary, not per component, so exit criterion~8
holds by construction: the logged command-application history of an
$L = \ell$ run is the $L = 0$ history delayed by exactly $\ell$ cycles,
with $\ell$ neutral cycles prepended. Both the computed and the applied
commands are logged, which is what makes the criterion checkable from
the log alone.

\subsection{Quaternion-error PD attitude controller}
\label{sec:gnc:pd}

Inputs per cycle: the commanded attitude and rate
$(\quat{q}_{i2\mathrm{cmd}},\, \vv{\omega}_{\mathrm{cmd}}^{\mathrm{CMD}})$
from the guidance component, and the estimated attitude and rate
$(\hat{\quat{q}}_{i2b},\, \hat{\vv{\omega}}^{B})$ from the navigation
component. The error quaternion is the attitude of the actual body
frame relative to the commanded frame, obtained from the composition
rule~\eqref{eq:notation:quatcomp}:
\begin{equation}
  \boxed{\;
  \delta\quat{q} = \quat{q}_{\mathrm{cmd2b}}
    = \quat{q}_{i2\mathrm{cmd}}^{*} \quatmul \hat{\quat{q}}_{i2b} ,
  \;}
  \label{eq:gnc:deltaq}
\end{equation}
so $\delta\quat{q} = [1, \vv{0}]$ at zero error, and for small errors
$\delta\vv{q}_v \approx \vv{\epsilon}/2$ with $\vv{\epsilon}$ the
rotation vector from commanded to actual body frame. The rate error
resolves the commanded rate into body axes through the error DCM,
\begin{equation}
  \vv{\omega}_{\mathrm{err}} = \hat{\vv{\omega}}^{B}
    - \dcm{\mathrm{CMD}}{B}(\delta\quat{q})\,
      \vv{\omega}_{\mathrm{cmd}}^{\mathrm{CMD}} ,
  \label{eq:gnc:werr}
\end{equation}
with the DCM from equation~\eqref{eq:notation:quat2dcm}. The commanded
torque, componentwise for $i \in \{x, y, z\}$ and with the shortest-path
sign
\begin{equation}
  s = \begin{cases} +1, & \delta q_w \ge 0,\\ -1, & \delta q_w < 0,
      \end{cases}
  \label{eq:gnc:sign}
\end{equation}
is
\begin{equation}
  \boxed{\;
  \tau_i = -\,K_{p,i}\; s\,\delta q_{v,i}
           \;-\; K_{d,i}\; \omega_{\mathrm{err},i} ,
  \;}
  \label{eq:gnc:pd}
\end{equation}
followed by the per-axis clamp
\begin{equation}
  \tau_i \leftarrow \min\bigl(\max(\tau_i,\, -\tau_{\max,i}),\,
    \tau_{\max,i}\bigr).
  \label{eq:gnc:sat}
\end{equation}
Normative arithmetic, for the exit-criterion-2 reproduction: the inputs
are used as received (no renormalization of $\delta\quat{q}$); the
Hamilton product~\eqref{eq:notation:hamiltonproduct} and conjugate are
evaluated in the project's scalar-first component order; the sign
branch is exactly $\delta q_w \ge 0$ (zero maps to $+1$); products and
clamps are elementwise with the gains $K_{p,i} > 0$
(\unit{\newton\metre}, per unit quaternion-vector component) and
$K_{d,i} > 0$ (\unit{\newton\metre\second\per\radian}). Because
$\delta\vv{q}_v \approx \vv{\epsilon}/2$, the small-error proportional
stiffness is $K_{p,i}/2$ per radian --- the factor of two lives in the
gain interpretation, not in the code. The $s$ factor selects the
shorter of the two rotations representing the same error (quaternion
double cover), which removes the classical unwinding behavior; global
asymptotic stability of the law for a rigid body with positive-definite
gains is the standard quaternion-feedback result
(Wie~\cite{wie2008}; Markley and Crassidis~\cite{markley2014}).

\subsection{Dead-reckoning attitude navigation}
\label{sec:gnc:deadreckoning}

The navigator holds a unit quaternion estimate
$\hat{\quat{q}}_{i2b}$, initialized from configuration
($\hat{\quat{q}}_0$; the reference scenarios initialize it to the true
initial attitude, which the mission file states explicitly --- there is
no implicit truth access). Each cycle it composes the measured gyro
increment $\Delta\tilde{\vv{\theta}}_k$
(equation~\eqref{eq:imu:gyro}) as an exact rotation. With the exact
exponential map
\begin{equation}
  \delta\quat{q}(\vv{\phi}) =
    \Bigl[\cos\tfrac{\phi}{2},\;
      \sin\tfrac{\phi}{2}\,\frac{\vv{\phi}}{\phi}\Bigr],
  \qquad \phi = \norm{\vv{\phi}},
  \qquad
  \delta\quat{q}(\vv{0}) = [1, \vv{0}],
  \label{eq:gnc:exactrot}
\end{equation}
the propagation and rate estimate are
\begin{equation}
  \boxed{\;
  \hat{\quat{q}}_{k} = \frac{\hat{\quat{q}}_{k-1} \quatmul
      \delta\quat{q}\bigl(\Delta\tilde{\vv{\theta}}_k\bigr)}
    {\bigl\lVert \hat{\quat{q}}_{k-1} \quatmul
      \delta\quat{q}\bigl(\Delta\tilde{\vv{\theta}}_k\bigr)
      \bigr\rVert},
  \qquad
  \hat{\vv{\omega}}_k^{B} =
    \frac{\Delta\tilde{\vv{\theta}}_k}{\Delta t_s} ,
  \;}
  \label{eq:gnc:drprop}
\end{equation}
body-side composition per equation~\eqref{eq:notation:quatcomp},
normalized every cycle. The exact-rotation choice (rather than a
small-angle quaternion) is pinned because it is branch-free above the
$\phi = 0$ special case, exactly norm-preserving before rounding, and
identical to the map the error-state filter of Chapter~\ref{ch:ekf}
uses --- one shared function. Treating the increment as a rotation
vector is exact for rotation about a fixed axis within the interval;
for axis motion within the interval the per-cycle residual is the
classical single-sample coning term, bounded by
\begin{equation}
  \norm{\vv{e}_{\mathrm{coning},k}}
    \le \tfrac{1}{12}\,
      \bigl\lVert \Delta\vv{\theta}_{k-1} \times \Delta\vv{\theta}_k
      \bigr\rVert
      + O\bigl(\norm{\Delta\vv{\theta}}^4\bigr),
  \label{eq:gnc:coningbound}
\end{equation}
the leading term of the two-sample coning correction this navigator
deliberately omits (Savage~\cite{savage1998a}; Titterton and
Weston~\cite{titterton2004}). At $\qty{100}{\hertz}$ cycling and
$\norm{\vv{\omega}} \le \qty{0.1}{\radian\per\second}$ the bound is
$\lesssim \qty{8e-8}{\radian}$ per second of flight --- far below the
gyro error budget, which is the point: the reference navigator's error
is dominated by the sensor, not the algorithm.

\subsection{Ascent pitch-program guidance component}
\label{sec:gnc:ascent}

The component is configured with the same data as the Phase~4 open-loop
sequence action: flight azimuth $A$, node tables $(t_j, \theta_j)$,
$j = 1 \dots P$, and it captures the frozen launch-site tangent basis
$(\hat{\vv{u}}_0, \hat{\vv{e}}_0, \hat{\vv{n}}_0)$ at initialization
(Chapter~\ref{ch:vehicle6dof}). Per cycle at time $t$:
\begin{enumerate}
  \item \textbf{Pitch interpolation, clamped:}
    \begin{equation}
      \theta(t) =
      \begin{cases}
        \theta_1, & t \le t_1,\\
        \theta_P, & t \ge t_P,\\
        \theta_j + \dfrac{t - t_j}{t_{j+1} - t_j}
          \bigl(\theta_{j+1} - \theta_j\bigr),
          & t_j < t \le t_{j+1} \text{ (first such } j\text{)},
      \end{cases}
      \label{eq:gnc:interp}
    \end{equation}
    the exact rule of the Phase~4 implementation, restated here so both
    paths interpolate identically.
  \item \textbf{Commanded attitude:} thrust axis
    $\hat{\vv{d}}^{I}(\theta(t), A)$ from
    equation~\eqref{eq:vehicle6dof:pitchaxis} in the frozen basis, and
    $\quat{q}_{i2\mathrm{cmd}}(t)$ from the Gram--Schmidt triad of
    equation~\eqref{eq:vehicle6dof:attitude} with up-reference
    $\hat{\vv{u}}_0$, including its degeneracy fallback (reference
    within $10^{-8}$ of parallel to the thrust axis: substitute
    $\hat{\vv{z}}$, or $\hat{\vv{x}}$ when $|\hat{d}_z| \ge 0.9$) ---
    the roll convention must match Phase~4 exactly for exit
    criterion~10's re-gate to compare like with like.
  \item \textbf{Commanded rate}, the finite-difference rotation between
    this cycle's and the next cycle's commanded attitude, resolved in
    the commanded frame at $t$:
    \begin{equation}
      \boxed{\;
      \delta\quat{q}_c = \quat{q}_{i2\mathrm{cmd}}(t)^{*} \quatmul
        \quat{q}_{i2\mathrm{cmd}}(t + T_c),
      \qquad
      \vv{\omega}_{\mathrm{cmd}}^{\mathrm{CMD}} =
        \frac{2\,\operatorname{sgn}(\delta q_{c,w})}{T_c}\,
        \delta\vv{q}_{c,v} ,
      \;}
      \label{eq:gnc:cmdrate}
    \end{equation}
    with $\operatorname{sgn}(0) = +1$: the small-angle rate extraction
    on the short branch of the double cover, exactly the Phase~4
    convention. Its $O(\norm{\delta\vv{q}_{c,v}}^3)$ small-angle error
    is negligible at per-cycle command increments
    ($\lesssim 10^{-2}$ rad).
\end{enumerate}
The component emits
$(\quat{q}_{i2\mathrm{cmd}}, \vv{\omega}_{\mathrm{cmd}}^{\mathrm{CMD}})$
for the PD controller of Section~\ref{sec:gnc:pd}; closing the loop
through the Chapter~\ref{ch:actuators} actuators is what exit
criterion~10 re-times against the FR-32 ascent performance target.

\section{Discretization and implementation}
\label{sec:gnc:implementation}

Binding implementation notes for the Phase~6 modules:
\begin{enumerate}
  \item \textbf{Modules and traceability.} Each component is a
    C++ \texttt{IGncComponent} registered in the FR-25 static registry
    and selected by configuration string; Python subclasses enter
    through the pybind11 trampoline with identical semantics. Source
    comments echo \texttt{eq:gnc:deltaq}, \texttt{eq:gnc:pd},
    \texttt{eq:gnc:sat}, \texttt{eq:gnc:werr},
    \texttt{eq:gnc:exactrot}, \texttt{eq:gnc:drprop}, and
    \texttt{eq:gnc:cmdrate} verbatim.
  \item \textbf{Statelessness inventory.} The PD controller is
    stateless; the dead reckoner's state is $\hat{\quat{q}}$; the
    guidance's state is the frozen basis and tables. All state is
    reset by \texttt{Sim.reset} (FR-24), which is what makes stepping
    and batch runs bit-identical (exit criterion~4).
  \item \textbf{No draws, no clock, no allocation.} Components consume
    no PRNG streams and never read wall-clock time; the update path
    allocates nothing (D-10).
  \item \textbf{Logged channels.} Computed and applied commands land in
    the reserved \texttt{nav.*} group; the dead reckoner logs
    $\hat{\quat{q}}$ (scalar-first) as its estimate channel. The EKF's
    richer \texttt{nav.est.*} layout is defined in
    Chapter~\ref{ch:ekf}.
  \item \textbf{Sensor-validity handling.} Components receive each
    sensor sample with its validity flag
    (Chapter~\ref{ch:sensors-optical} semantics) and must skip work
    keyed to invalid samples; the dead reckoner consumes only the IMU,
    which is always valid.
\end{enumerate}

\section{Validation evidence}
\label{sec:gnc:validation}

This chapter precedes its modules; the validation-evidence table with
machine-checked test identifiers is added by the implementing change.
The planned gates, mirroring Phase~6 exit criteria~2, 4, 8, and~10:
\begin{itemize}
  \item \textbf{Python PD reproduction (exit criterion 2).} A pure
    Python transcription of
    equations~\eqref{eq:gnc:deltaq}--\eqref{eq:gnc:sat}, written from
    this chapter alone, reproduces the C++ controller's commanded
    torques to $< \qty{e-9}{\newton\metre}$ at every cycle of a golden
    scenario that exercises both signs of $\delta q_w$, all clamp
    branches, and nonzero commanded rates.
  \item \textbf{PD regulation properties.} From a large-angle initial
    error under torque-only rigid-body dynamics
    (Chapter~\ref{ch:rigidbody}), the closed loop converges to the
    commanded attitude with the expected linearized time constants;
    initializing at $\delta q_w < 0$ takes the short path (no
    unwinding).
  \item \textbf{Dead-reckoning drift.} On the noise-free coning
    scenario of Chapter~\ref{ch:rigidbody}, the dead reckoner's error
    stays within the accumulated bound of
    equation~\eqref{eq:gnc:coningbound}; with a biased gyro the drift
    matches the injected bias integral to first order.
  \item \textbf{Guidance equivalence.} On the Phase~4 reference ascent,
    the component's commanded attitude history equals the Phase~4
    open-loop history to rounding (same interpolation, same triad,
    same fallback), and the closed-loop trajectory with the PD
    controller and actuators meets its dispersion budget against the
    open-loop baseline.
  \item \textbf{Latency shift (exit criterion 8).} For $L \in \{1, 5\}$
    the logged applied-command history equals the $L = 0$ history
    shifted by exactly $L$ cycles with neutral commands prepended,
    checked channel-by-channel.
  \item \textbf{Performance re-gate (exit criterion 10).} The FR-32
    ascent target ($\ge 100\times$ real time, Pi~5 single core, per
    the Phase~5 deferral provisions) holds with the full built-in GNC
    stack in the loop.
\end{itemize}

\section{References}
\label{sec:gnc:references}

The quaternion-feedback PD law and its stability properties follow
Wie~\cite{wie2008} and Markley and Crassidis~\cite{markley2014},
transcribed to the project's D-7 conventions; the strapdown increment
composition and the coning-residual bound follow
Savage~\cite{savage1998a} and Titterton and
Weston~\cite{titterton2004}; the ascent geometry is derived in
Chapter~\ref{ch:vehicle6dof}. Full bibliographic details appear in the
bibliography.
